\section{WHIR}
\label{s:WHIR}


% Given the commitment $g_0$ of a multil

% Given a function $g_0\in F^{D}$, and proximity parameter $\theta$, WHIR \cite{WHIR} proves proximity of $g_0$ to a polynomial $u_0(X)$ the multilinear representation of which,  $P_0(X_1,\ldots, X_{k_0})$, satisfies a linear claim of the form
% \[
% c_0 = \langle P_0, K_0\rangle_{H_{k_0}},
% \]

% Given an inner product claim $c_0 = \langle P_0, K_0\rangle_{H_{k_0}}$ between a multilinear $P_0(X_1,\ldots, X_m)$ with a public $K_0(X_1,\ldots, X_n)$, WHIR recursively transforms it into smaller claims 
% \[
% c_i = \langle P_i(\:.\:), K_i(\:.\:)\rangle_{H_{m-k_i}},
% \]
% where $P_i$ and $K_i$ are random linear combinations of the previous $P_{i-1}$ and $K_{i-1}$, respectively. 

% is a random specialization of the previous polynomial $P_{i-1}$, and $K_i$ an aggregate of the previous $K_{i-1}$ together with the their consistency spot checks.

For completeness we give a concise description of the WHIR \cite{WHIR} proof of proximity to constrained Reed--Solomon codes. 
Given a word defined on some set, $g_0:D\longrightarrow \mathbb F_q$, WHIR \cite{WHIR} is an interactive oracle proof of proximity to a univariate polynomial $u_P(X)$, the multivariate representation of which, say $P(X_1, \ldots, X_{k_0})$, satisfies a linear claim of the form
\begin{equation}
\label{e:WHIR:main:claim}
c_0 = \langle P_0, K_0\rangle_{H_{k_0}},
\end{equation}
where $K_0(X_1,\ldots, X_{k_0})$ is a publicly known (and often succinctly evaluable) multilinear. 
Similar to Basefold \cite{Basefold} and preceding works \cite{Gemini,Hyperplonk} \cite{DualGemini:hackmd,Zeromorph}, the construction rests on the connection between multilinear sumcheck specialization and the FRI-like folding of their univariate representations.

\subsection{Setup}
The mechanics of WHIR is in alignment with a specified sequence of univariate spaces
\begin{align}
\label{e:chain:function:spaces}
\mathscr P_{k_0} \supset \mathscr P_{k_1} \supset \ldots \supset \mathscr P_{k_r},
\end{align}
with $k_{0} > k_1 > \ldots > k_{r}\geq 0$. 
(These will correspond to the number of remaining variables after each round.)
As before, we commit polynomials in interleaved representation, using the linear isomorphism
\[
\mathscr P_{k_i} \simeq \mathscr P_{k_{i+1}}^{e_{i+1}}, \quad e_{i+1} = \dim{\mathscr P_{k_i}/\mathscr P{k_{i+1}}},
\]
by FFT decomposition \eqref{e:FFT:decomp}, for each $i=0,\ldots, r-1$.
The decomposed polynomials, elements from
\begin{equation}
\mathscr P_{k_1}^{e_1}, \mathscr P_{k_2}^{e_2}, \ldots, \mathscr P_{k_{r}}^{e_{r}},
\end{equation}
are then %evaluated over their own domain $D_i$, 
committed as words from the interleaved codes
\begin{equation}
 C_1^{e_1}, C_2^{e_2}, \ldots, C_r^{e_r},
\end{equation}
with specified base codes $C_{i} = \RS\left[F, D_{i}, 2^{k_{i}}\right]$ of (typically increasing) minimum distance.
Their evaluation domains $D_i$ are chosen as multiplicative cosets\footnotemark in $\mathbb F_q$, and are of size $|D_{i}|= 2^{n_{i}}$ for some $n_i > k_{i}$. 
\footnotetext{%
More generally, disjoint unions of multiplicative cosets of size $2^{k_i}$.
}
Each round further specifies $s_i\geq 1$ for the number of samples in the consistency check between adjacent oracles, as described below.


\subsection{The IOPP.}

%In the following $\mathbb F_q$ is an FFT-friendly prime field, and $F$ a finite extension of it. 
The protocol starts with initial claim \eqref{e:WHIR:main:claim} and consists of $r$ rounds. 
Each round $i=1,\ldots,r$ reduces a claim $c_{i-1} = \langle P_{i-1}, K_{i-1}\rangle_{H_{k_{i-1}}}$ on the previous polynomial $P_{i-1}$ to a claim $c_{i} = \langle P_{i}, K_{i}\rangle_{H_{k_i}}$ on a new polynomial.

%typically cryptographically large.

% As described at the beginning of Section \ref{s:protocol:no:zk}, WHIR specifies a chain of univariate polynomial spaces
% \begin{align*}
% % \label{e:chain:function:spaces}
% \mathscr P_{k_0} \supset \mathscr P_{k_1} \supset \ldots \supset \mathscr P_{k_{r-1}} \supset \mathscr P_{k_r},
% \end{align*}
% with $k_{0} > k_1 > \ldots > k_{r}\geq 0$.
% A polynomial $p(X)$ from $\mathscr P_{k_{i-1}}$, $i = 1,\ldots, r$, is committed in decomposed manner as element $(p_1(X),\ldots, p_{e_i}(X))$ from  $\mathscr P_{k_i}^{e_i}$, where $e_i = \dim \mathscr P_{k_{i-1}} / \mathscr P_{k_i}$, according to the common FFT decomposition \eqref{e:FFT:decomp}.
% Thus, the codes met in the course of WHIR are the interleaved Reed--Solomon codes $C_{i}^{e_i}$, $i=1,\ldots, r$, with base code
% \[
% C_{i} = \RS[F,D_i,2^{k_i}], \quad i=1,\ldots, r,
% \]
% generated by the polynomials from $\mathscr P_{k_i}$, respectively.
% The evaluation domains $D_i$ are multiplicative subgroups of $\mathbb F_q$, of size $|D_i|=2^{n_i}$, for some $n_i > k_i$.


\begin{protocol}[WHIR]
\label{prot:WHIR}
The prover has a multilinear $P_{0}\in \mathbb F_q[X_1,\ldots, X_{k_0}]$, satisfying the inner-product claim
\[
c_0 = \langle P_{0}, K_0 \rangle_{H_{k_0}} 
\]
with some public multilinear $K_{0}\in F[X_1,\ldots, X_{k_0}]$.
The verifier has oracle access to $g_{0}\in (\mathbb F_q^D)^{e_1}$ the commitment of $P_0$'s univariate representation $u_{P_0}(X)\in \mathscr P_{k_0}$ as interleaved word from $C_1^{e_1}$.

Both prover and verifier engage in the following $r$ rounds, for $i=1,\ldots, r$, each consisting of three phases. 
\begin{enumerate}
    \item 
    Sumcheck for $r_i = k_{i-1} - k_i$ rounds, on the previous inner product claim between $P_{i-1}$ and $K_{i-1}$.
    In the course of these rounds, the prover has sent the quadratic refinements polynomials $q_{i,1}(X)$, \ldots $q_{i,r_i}(X)$, and the verifier has answered with $\vec\beta_i = (\beta_{i,1}, \ldots \beta_{1, r_i})\in F^{r_i}$.
    The prover commits to 
    \[
    P_{i} = P_{i-1}(\vec \beta_{i},\:.\:), 
    \]
    % the latter only when $k_i +1 \leq m_2$, 
    as the polynomial $u_{P_i}(X)\in \mathscr P_{k_i}$ represented by the interleaved word
    \[
    g_{i}  \in C_{i+1}^{e_{i+1}},
    \]
    and provides the verifier oracle access.
    (In the case $i=r$, the prover reveals the polynomial $u_{P_r}(X)$ in plain.)
    % and their next refinement polynomial
    % \begin{align*}
    % q_{k_i + 1}(X) = &\langle K_{i,1}(\:.\:,X,\beta_{k_i + 1}, \ldots,\beta_{m_1}), p_1(\:.\:,X,\beta_{k_i + 1}, \ldots, \beta_{m_1})\rangle
    % \\
    % &+ \langle K_2(\:.\:,X,\beta_{k_i + 1}, \ldots,\beta_{m_2}), 
    %  p_2(\:.\:,X,\beta_{k_i + 1}, \ldots,\beta_{m_2}) \rangle,
    % \end{align*}
    \item 
    Query phase. 
    \begin{enumerate}
    \item
    The verifier draws a DEEP query $x_{i,0} \sample F$, on which the prover answers with the value $v_{0}$ of the univariate representations $u_{P_i}(X)$ at $x_{0}$.  
    % (In the case $i=r$, this step can be dropped.)
    \item 
    The verifier queries $g_{i-1}$ from the previous round at $s_i>0$ random positions 
    \[
    x_{i,1}, \ldots, x_{i,s_i} \sample D_{i},
    \]
    and computes the expected values $v_{1}, \ldots, v_{s_i}$ for $u_{P_{i}}(X)$ at these points from the oracle answers, and $\vec\beta_i$.
    \end{enumerate}
    \item 
    Aggregation.
    The verifier samples $\gamma_i\sample F$ for combining the specialized sumcheck claim
    \begin{align*}
    q_{i, r_i}(\beta_{i,r_i}) &= \langle P_{i}, K_{i-1}(\vec \beta_{i},\:.\:) \rangle_{H_{k_i}} 
    \end{align*}
    and the evaluation claims
    \begin{align*}
    v_{j} &= \langle P_{i}, L_{x_i}\rangle_{H_{k_i}}, 
    \quad j = 0,\ldots, s_i,
    \end{align*}
    into a the new inner product claim
    $
    c_i = \langle P_{i}, K_{i}\rangle_{H_{k_i}} 
    $
    with the new $c_i$ and kernel $K_{i}$ as clarified below.
\end{enumerate}


If the consistency checks in all the sumcheck rounds hold, and if the final inner product claim $c_r = \langle P_r, K_r\rangle$ holds for the multilinear representation of the revealed polynomial $u_{P_r}(X)$, then the verifier accepts.
\end{protocol}

In the aggregation step, the combined kernel is
\begin{align*}
    K_{i} &=  K_{i-1,i,1}(\vec\beta_{i},X_1,\ldots, X_{k_i}) + \sum_{j=0}^{s_i} \gamma_i^{j+1} \cdot L_{x_{j}}(X_1,\ldots, X_{k_i}),
    \end{align*}
and $c_i = q_{i,r_i}(\beta_{i,r_i}) + \sum_{j=0}^{s_i} \gamma_i^{j+1}\cdot v_{j}$.
    

% \subsection{Combining several claims.}



\subsection{Soundness.}
\label{s:WHIR:soudness:discussion}

% Let us quickle discuss the soundness properties of WHIR, for proximity parameters in the Johnson regime of the Reed--Solomon codes, 
% \[
% \theta_i \in \left(0, 1 -\sqrt{1-\delta_i}\right),
% \]
% where $\delta_i$ is the minimum distance of $C_i$, $i=1,\ldots, r$.

Every round of WHIR, $i=1,\ldots, r$, is a randomized reduction $\mathcal R_{i-1}\longrightarrow \mathcal R_i$ in the verifier randomnesses $\beta_{i,1},\ldots,\beta_{i,r_i}$, $x_{i,0},x_{i,1},\ldots, x_{i,s_i}$, and $\gamma_i$, where 
\begin{equation}
\label{e:WHIR:round:relation}
\mathcal R_i = \left\{
\begin{matrix}
g_i\in (\mathbb F_q^{D_{i+1}})^{e_{i+1}}
% \\
% q_0(X) \in F[X]^{\leq 2}
\end{matrix}
\::\:
\begin{matrix}
\exists p_i(X)\in \mathscr P_{k_{i+1}}^{e_{i+1}} \hat = \mathscr P_{k_i}, 
\\
d(g_i,p_i)\leq \theta_{i+1},
\\
\langle P_i, K_i\rangle_{H_{k_i}} = c_i 
% \\
% q_0(0) + q_0(1) = c_0
\end{matrix}
\right\},
\end{equation} 
for $i=0,\ldots, r-1$, again $P_i$ is the multilinear representation of $p_i$ as polynomial from $\mathscr P_{k_i}$, and the distance refers to the interleaved representation of $p_i$ as codeword from $C_{i+1}^{e_{i+1}}$. 
In the last case $i=r$,
\begin{equation}
\label{e:WHIR:last:relation}
\mathcal R_r = \big\{
p_r(X)\in \mathscr P_{k_r} 
\::\:
\langle P_r, K_r\rangle_{H_{k_r}} = c_r 
\big\}.
\end{equation} 
For brevity we use lazy notation here, hiding the dependency of $\mathcal R_i$ on the previous messages:
The inner product kernel $K_i$ incorporates the specialized sumcheck claim on the previous oracle $g_{i-1}$, its values at the $s_i$ in-domain queries, and the claimed value for $p_i$ at the single out-of-domain point $x_{i,0}$.
% We note that in the case $i=r$, since $p_r(X)\in \mathscr P_{k_r}$ is revealed in plain, $g_r = p_r$, the distance condition in \eqref{e:WHIR:round:relation} becomes vacuous (besides that $\theta_{r+1}$ is not specified).

% The soundness analysis of the sumcheck phase, in parallel with the multilinear combinations of the interleaved components, is the most technically involved part. 
% Its proof relies on the curve decodability of Reed--Solomon codes, yielding correlated agreement under the additional constraints imposed by the inner products. 
% for \textit{constrained} Reed--Solomon codes \cite[Theorem 4]{Basefold:LDR}, a consequence of degree-2 curve decodability up to Johnson radius of $C_i$, 
% %which is a consequence of the more recent notion of line decodability \cite{goyalguruswami}.
% which we shall state in a simplified form in Appendix \ref{s:CAT}.
%
Without specifying the intermediate relations of each step of a WHIR round, we state its soundness properties as an (interactive oracle) reduction.
For a detailed analysis we refer to \cite[Section 6.4]{Basefold:LDR}.

\begin{thm}[WHIR Soundness]
\label{thm:WHIR:soundness}
For $\theta_i\in (0,1 -\sqrt{1-\delta_i})$, where $\delta_i$ is the minimum distance of the Reed--Solomon code $C_i$.
Then round $i=1,\ldots, r-1$ in Protocol \ref{prot:WHIR} is round-by-round (knowledge) sound oracle reduction $\mathcal R_{i-1}\longrightarrow \mathcal R_{i}$, with the following soundness errors for its steps.
\begin{enumerate}
    \item 
    Sumcheck phase:
    For each of the $r_i$ steps,
    \begin{align*}
       \varepsilon_{S,i} &= \varepsilon_{\text{CA}}(C_{i},2,\theta_{i}), 
    \end{align*}
    where $\varepsilon_\text{CA}(C_{i},2,\theta_{i})$ is the soundness error of the constrained correlated agreement theorem (Theorem \ref{thm:CAT:constrained}) for $C_{i}$, under degree-$2$ constraints, using proximity parameter $\theta_{i}$.
    \item 
    Query phase:
    \begin{align*}
        \varepsilon_{Q,i} = \frac{\ell_{i+1}(\theta_{i+1})^2}{2} \cdot \frac{2^{k_{i+1}}}{|F\setminus D_{i+1}|} + (1 -\theta_{i})^{s_i},
    \end{align*}
    where $\ell_{i+1}(\theta_{i+1})$ is the list size bound for $C_{i+1}$ at distance $\theta_{i+1}$.
    \item 
    Aggregation step:
    \begin{align*}
        \varepsilon_{A,i} = \ell_{i+1}(\theta_{i+1}) \cdot \frac{s_{i}}{|F|},
    \end{align*}
    again with the same list size bound as above.
\end{enumerate}
The same assertion is true for last round $i=r$, except that the query phase error $\varepsilon_{Q,r} = (1 - \theta_r)^{s_r}$, and the aggregation step error $\varepsilon_{A,r} = \frac{s_i}{|F|}$.
\end{thm}



% \begin{rem}
% For the soundness analysis, one may alternatively use the high-level notion of mutual correlated agreement \cite{WHIR}, a consequence of curve-decodability \cite{cryptoeprint:2025/2110,cryptoeprint:2025/2054}. 
% \end{rem}
% Let us roughly discuss the soundness errors. 
% \begin{rem}
% Alternatively, one may use the stronger notion of mutual correlated agreement.
% \end{rem}

% Each step of the sumcheck is a randomized reduction $\mathcal R_{i-1}\stackrel{\lambda_i}{\longrightarrow} \mathcal R_i$, where
% \[
% \mathcal R_i = \left\{
% \begin{matrix}
% f_0\in F^{D_0},
% \\
% q_0, \ldots, q_{i-1} \in F[X]^{\leq 2}
% \end{matrix}
% \::\:
% \begin{matrix}
% q_{j-1}(\lambda_j) = q_j(0) + q_j(1), j=1,\ldots,i-1,
% \\
% \text{and}
% \\
% \exists p_i(X)\in \mathscr P_{n-i}, %P_i =\Phi^{-1}(p_i),
% \\
% d_{\pi^i(D_0)}(\mathsf{fold}_{\lambda_1,\ldots, \lambda_i}(f_0),p_i)\leq \theta_0,
% % \\\text{s.t.} 
% \\
% \langle P_i, R_0(\lambda_1,\ldots, \lambda_i,\:.\:)\rangle = q_{i-1}(\lambda_i),
% \end{matrix}
% \right\},
% \]
% for $i=1,\ldots, k-1$.